% !TeX root=main.tex
\chapter{نتیجه‌گیری و کارهای آینده}
\section{جمع‌بندی پژوهش}

در این پایان‌نامه، مسئله مقیاس‌پذیری بلاکچین با تمرکز بر سامانه‌های مبتنی بر مدل خروجی تراکنش خرج‌نشده\LTRfootnote{Unspent Transaction Output (UTXO)} بررسی شد. ابتدا نشان داده شد که بیتکوین با وجود امنیت و غیرمتمرکز بودن بالا، به دلیل معماری تک‌زنجیره‌ای، فاصله بلوک طولانی و محدودیت اندازه/وزن بلوک، توان عملیاتی محدودی دارد. سپس راهکارهای مختلف مقیاس‌پذیری شامل افزایش ظرفیت لایه اول، کانال‌های پرداخت، رول‌آپ‌ها، شاردینگ، زنجیره‌های جانبی و تغییر الگوریتم اجماع بررسی شدند.

در ادامه، معماری پیشنهادی بر پایه زنجیره اصلی و زنجیره‌های جانبی موازی معرفی شد. در این معماری، زنجیره اصلی نقش لنگر امنیتی و محل ثبت ارجاعات را دارد و زنجیره‌های جانبی تراکنش‌ها را بر اساس تابع تخصیص قطعی پردازش می‌کنند. مدل ریاضی نشان داد که اگر ظرفیت هر زنجیره جانبی ثابت باشد، توان عملیاتی کل تقریباً به صورت خطی با تعداد زنجیره‌ها رشد می‌کند.

\section{پاسخ به پرسش‌های پژوهش}

پرسش نخست درباره تقسیم تراکنش‌ها بدون افزایش غیرقابل کنترل دوبارخرج بود. پاسخ پژوهش، استفاده از تابع تخصیص مبتنی بر هش خروجی‌های خرج‌نشده است که هر خروجی را به قلمرو مشخصی از زنجیره جانبی نگاشت می‌کند. پرسش دوم درباره نقش زنجیره اصلی بود. در معماری پیشنهادی، زنجیره اصلی محل پردازش همه تراکنش‌ها نیست، بلکه لنگر امنیتی، زمان‌سنج و مرجع نهایی ارجاعات است. پرسش سوم درباره رابطه تعداد زنجیره‌های جانبی و توان عملیاتی بود. مدل نشان داد تحت فرض‌های تقارنی، توان کل با \(N\) رشد خطی دارد. پرسش چهارم درباره تفاوت توان عملیاتی نظری و بنچمارک\LTRfootnote{Benchmark} بود. نتایج نشان داد بنچمارک محاسباتی نباید با توان واقعی شبکه اشتباه گرفته شود. پرسش پنجم درباره ریسک‌ها بود. ریسک‌های اصلی شامل امنیت زنجیره جانبی، تراکنش بین‌زنجیره‌ای، داده‌دسترسی، هماهنگی و عدم تعادل بار هستند.

\section{دستاوردهای اصلی}

دستاوردهای اصلی پایان‌نامه عبارت‌اند از:

\begin{itemize}
\item ارائه یک تکسونومی تحلیلی از روش‌های مقیاس‌پذیری بلاکچین.
\item طراحی معماری چندزنجیره‌ای مبتنی بر زنجیره اصلی و زنجیره‌های جانبی.
\item تعریف تابع تخصیص قطعی بر اساس خروجی‌های خرج‌نشده.
\item ارائه مدل ریاضی برای توان عملیاتی و رشد خطی ظرفیت.
\item پیاده‌سازی نمونه اولیه در زبان راست\LTRfootnote{Rust}.
\item تفکیک صریح میان بنچمارک محاسباتی و ظرفیت شبکه واقعی.
\item شناسایی چالش‌های امنیتی و مهندسی برای توسعه آینده.
\end{itemize}

\section{محدودیت‌ها}

این پژوهش محدودیت‌هایی دارد. نخست، پیاده‌سازی نمونه اولیه همه جزئیات شبکه عمومی را مدل نمی‌کند. دوم، پروتکل تراکنش بین‌زنجیره‌ای در نسخه اولیه نیازمند تکمیل و اثبات امنیتی دقیق‌تر است. سوم، امنیت زنجیره‌های جانبی با دشواری پایین‌تر باید با تحلیل حمله و مدل اقتصادی تکمیل شود. چهارم، داده‌دسترسی بلوک‌های جانبی و سیاست ذخیره‌سازی نودها باید دقیق‌تر طراحی شود. پنجم، نتایج بنچمارک به سخت‌افزار و تنظیمات آزمایش وابسته‌اند.

\section{کارهای آینده}

برای توسعه این پژوهش، مسیرهای آینده زیر پیشنهاد می‌شود:

\begin{enumerate}
\item طراحی و پیاده‌سازی پروتکل کامل بین‌زنجیره‌ای با قفل اتمیک\LTRfootnote{Atomic Lock} و اثبات مرکل.
\item بررسی استفاده از استخراج ادغام‌شده\LTRfootnote{Merged Mining} برای افزایش امنیت زنجیره‌های جانبی.
\item طراحی سیاست دوره‌بندی\LTRfootnote{Epoch} برای تغییر تعداد زنجیره‌های جانبی بدون تغییر ناگهانی تخصیص خروجی‌ها.
\item افزودن داده‌دسترسی و نمونه‌برداری داده برای بلوک‌های جانبی.
\item اجرای بنچمارک توزیع‌شده روی چند ماشین و شبیه‌سازی تأخیر شبکه.
\item تحلیل اقتصادی کارمزدها، پاداش ماینرها و انگیزه مشارکت در زنجیره‌های جانبی.
\item مقایسه تجربی با پیاده‌سازی‌های مشابه مانند زنجیره جانبی، شاردینگ و رول‌آپ در محیط کنترل‌شده.
\item وارسی رسمی\LTRfootnote{Formal Verification} پروتکل نهایی‌سازی و تراکنش بین‌زنجیره‌ای.
\end{enumerate}

\section{نتیجه نهایی}

معماری پیشنهادی این پایان‌نامه نشان می‌دهد که با جداسازی نقش امنیت و پردازش، می‌توان به مدلی دست یافت که نسبت به بلاکچین تک‌زنجیره‌ای ظرفیت بالاتری ارائه دهد و همچنان از زنجیره اصلی امن برای قطعیت نهایی استفاده کند. این معماری جایگزین نهایی همه راهکارهای مقیاس‌پذیری نیست، اما چارچوبی قابل تحلیل برای ترکیب مزایای زنجیره اصلی اثبات کار\LTRfootnote{Proof of Work (PoW)}، زنجیره‌های جانبی و پردازش موازی فراهم می‌کند. نتیجه اصلی پژوهش آن است که مقیاس‌پذیری واقعی نه از یک ترفند منفرد، بلکه از طراحی دقیق رابطه میان اجماع، داده، اعتبارسنجی، امنیت و تجربه کاربری حاصل می‌شود.

\section{دستاورد روش‌شناختی پژوهش}

در کنار خود معماری، رویکرد گزارش‌دهی این پژوهش نیز بخشی از سهم آن است. در بخش قابل توجهی از کارهای فنی این حوزه، عددهای عملکردی بدون زمینه کافی گزارش می‌شوند و خواننده نمی‌تواند به پرسش‌های اساسی پاسخ دهد: آیا محیط آزمایش شباهتی به شبکه واقعی داشته است؟ آیا محدودیت اندازه بلوک اعمال شده است؟ آیا امضای دیجیتال واقعی اعتبارسنجی شده یا از آن صرف‌نظر شده است؟ آیا تأخیر انتشار شبکه لحاظ شده است؟ آیا نقطه مقایسه منصفانه انتخاب شده است؟ عددی که نتوان این پرسش‌ها را درباره‌اش پاسخ داد، هر قدر هم بزرگ باشد، قابل اتکا نیست.

در این پژوهش، این پرسش‌ها از ابتدا در ساختار ارزیابی تعبیه شدند: پارامترهای آزمایش جداگانه جدول‌بندی شد، عددهای بزرگ پیاده‌سازی به زبان راست صراحتاً «ظرفیت محاسباتی محیط آزمایش» نامیده شدند نه توان عملیاتی شبکه عمومی، تهدیدهای اعتبار داخلی، خارجی، سازه‌ای و بازتولید یک‌به‌یک فهرست شدند، و مواردی که فقط در حد طراحی مفهومی باقی مانده‌اند از موارد پیاده‌سازی‌شده تفکیک شدند. این شفافیت شاید عددهای گزارش‌شده را کم‌زرق‌وبرق‌تر کند، اما ادعاهای پژوهش را قابل دفاع و قابل بازتولید نگه می‌دارد و مسیر کارهای بعدی را دقیق‌تر مشخص می‌کند.

\section{مسیر تبدیل پایان‌نامه به مقاله}

این پایان‌نامه می‌تواند مبنای یک مقاله علمی نیز باشد. برای تبدیل آن به مقاله، پیشنهاد می‌شود تمرکز مقاله بر سه بخش قرار گیرد: مدل رسمی معماری، تابع تخصیص مبتنی بر خروجی‌های خرج‌نشده و ارزیابی مقیاس‌پذیری. مقاله باید کوتاه‌تر از پایان‌نامه باشد، اما آزمایش‌ها باید دقیق‌تر شوند. به طور خاص، اجرای شبکه‌ای با چند نود فیزیکی یا کانتینر، گزارش تأخیر و نرخ فورک\LTRfootnote{Fork Rate}، و مقایسه با خطوط مبنای\LTRfootnote{Baseline} ساده مانند زنجیره منفرد سریع می‌تواند مقاله را قوی‌تر کند.

ساختار پیشنهادی مقاله چنین است: مقدمه و مسئله، کارهای مرتبط، معماری پیشنهادی، تحلیل امنیت، ارزیابی، بحث محدودیت‌ها و نتیجه‌گیری. شکل تکسونومی می‌تواند به عنوان شکل مرور ادبیات در مقاله استفاده شود، اما شکل معماری و نمودار رشد توان عملیاتی مهم‌تر خواهند بود.
